```latex
\documentclass[11pt]{amsart}

\usepackage{amsmath,amssymb,amsthm,mathtools}
\usepackage[colorlinks=true,linkcolor=blue,citecolor=blue,urlcolor=blue]{hyperref}

\numberwithin{equation}{section}
\newtheorem{theorem}[equation]{Theorem}
\newtheorem{proposition}[equation]{Proposition}
\newtheorem{lemma}[equation]{Lemma}
\newtheorem{corollary}[equation]{Corollary}
\theoremstyle{definition}
\newtheorem{definition}[equation]{Definition}
\theoremstyle{remark}
\newtheorem{remark}[equation]{Remark}

\newcommand{\kk}{\mathbb{k}}
\newcommand{\Uc}{\mathcal{U}}
\newcommand{\bcS}{\mathcal{S}}
\newcommand{\fB}{\mathfrak{B}}
\newcommand{\fb}{\mathfrak{b}}
\newcommand{\sln}{\mathfrak{sl}(n)}
\newcommand{\bipr}{\mathbin{\bowtie}}
\newcommand{\la}{\langle}
\newcommand{\ra}{\rangle}
\newcommand{\id}{\mathrm{id}}

% Web sources used below:
% [staff.fnwi.uva.nl](https://staff.fnwi.uva.nl/j.v.stokman/quantumdouble2.pdf)
% [staff.fnwi.uva.nl](https://staff.fnwi.uva.nl/j.v.stokman/Quantumsl2.pdf)
% [arxiv.org](https://arxiv.org/pdf/math/0512563)
% [cambridge.org](https://www.cambridge.org/core/services/aop-cambridge-core/content/view/886B0B4D8E3D242A51A0380A486351EA/S0013091504000689a.pdf/presentation-by-borel-subalgebras-and-chevalley-generators-for-quantum-enveloping-algebras.pdf)

\title[The punctured bigon and a quantum double]{The stated skein algebra of the
punctured bigon as a generalized quantum double}
\date{}

\begin{document}

\maketitle

\section{Conventions}

Let \(\kk=\mathbb{C}(q)\), where \(q\) is transcendental.  All results below remain
valid after specialization to an invertible \(q\in\mathbb{C}^{\times}\) which is not a
root of unity.  Put
\[
\lambda:=q^{-1}-q.
\]

We use the simply connected, or extended, quantum group
\(\Uc_q^{\mathrm{ext}}(\sln)\) in the conventions of Klimyk and Schm\"udgen
\cite{KS}.  Thus it has generators
\[
\widehat K_i^{\pm1}\quad(1\le i\le n),\qquad
E_i,\ F_i\quad(1\le i<n),
\]
subject to the following relations:
\begin{align}
\widehat K_i\widehat K_i^{-1}
&=\widehat K_i^{-1}\widehat K_i=1, &&
[\widehat K_i,\widehat K_j]=0,
\label{eq:Kcomm}\\
\prod_{i=1}^{n}\widehat K_i&=1, &&
\label{eq:Kprod}\\
\widehat K_iE_j\widehat K_i^{-1}
 &=q^{\delta_{ij}-\delta_{i,j+1}}E_j, &&
\widehat K_iF_j\widehat K_i^{-1}
 =q^{-\delta_{ij}+\delta_{i,j+1}}F_j,
 \label{eq:Kaction}\\
[E_i,F_j]
 &=\delta_{ij}\frac{K_i-K_i^{-1}}{q-q^{-1}},
 &&K_i:=\widehat K_i\widehat K_{i+1}^{-1},
 \label{eq:EFrel}
\end{align}
together with the quantum Serre relations
\begin{equation}\label{eq:ESerre}
E_i^2E_j-(q+q^{-1})E_iE_jE_i+E_jE_i^2=0
\qquad (|i-j|=1),
\end{equation}
\begin{equation}\label{eq:FSerre}
F_i^2F_j-(q+q^{-1})F_iF_jF_i+F_jF_i^2=0
\qquad (|i-j|=1),
\end{equation}
and
\[
[E_i,E_j]=[F_i,F_j]=0\qquad (|i-j|>1).
\]

The Borel subalgebras are
\[
\Uc_q(\fb_+)
 :=\kk\langle \widehat K_i^{\pm1},E_j\rangle,
\qquad
\Uc_q(\fb_-)
 :=\kk\langle \widehat K_i^{\pm1},F_j\rangle .
\]
Their standard skew pairing is the one in the question:
\begin{equation}\label{eq:pairing}
\la \widehat K_i,\widehat K_j\ra
   =q^{\frac1n-\delta_{ij}},\qquad
\la E_i,F_j\ra
   =\frac{\delta_{ij}}{q^{-1}-q},
\qquad
\la \widehat K_i,F_j\ra
   =\la E_j,\widehat K_i\ra=0.
\end{equation}
The product relation \eqref{eq:Kprod} is compatible with
\eqref{eq:pairing}, since
\[
\sum_{i,j}\left(\frac1n-\delta_{ij}\right)=0.
\]

The notation \(\bcS_n(\fB^*)\) denotes the stated \(\mathrm{SL}_n\)-skein algebra
of the once-punctured bigon, with the usual normalization in which the two
families of horizontal right-directed arcs are denoted \(b_{ij}\) and \(c_{ij}\).
No additional localization will be introduced: invertibility of the diagonal
arcs occurring below is part of the standard punctured-bigon presentation.

\section{The generalized quantum double}

It is useful to distinguish the two copies of the Cartan generators in the
double.  Write
\[
\widehat K_i^{+}\in \Uc_q(\fb_+),
\qquad
\widehat K_i^{-}\in \Uc_q(\fb_-),
\]
and let
\[
K_i^{+}
 :=\widehat K_i^{+}(\widehat K_{i+1}^{+})^{-1},
\qquad
K_i^{-}
 :=\widehat K_i^{-}(\widehat K_{i+1}^{-})^{-1}.
\]

\begin{proposition}[Presentation of the generalized double]
\label{prop:double-presentation}
Let
\[
D_q:=\Uc_q(\fb_-)\bipr\Uc_q(\fb_+)
\]
be the generalized quantum double determined by the skew pairing
\eqref{eq:pairing}.  As an algebra, \(D_q\) is generated by
\[
\widehat K_i^{+\pm1},\quad
\widehat K_i^{-\pm1},\quad
E_i,\quad F_i,
\]
subject to the following relations.

\begin{enumerate}
\item Each sign separately satisfies the Borel relations
\eqref{eq:Kcomm}--\eqref{eq:ESerre} or
\eqref{eq:Kcomm}--\eqref{eq:Kaction} and \eqref{eq:FSerre}.

\item The two Cartan tori commute:
\begin{equation}\label{eq:D-Cartan-commute}
[\widehat K_i^{+},\widehat K_j^{-}]=0.
\end{equation}

\item Both Cartan tori act diagonally on both nilpotent parts:
\begin{align}
\widehat K_r^{\eta}E_i
 &=q^{\delta_{ri}-\delta_{r,i+1}}
   E_i\widehat K_r^{\eta},
 \label{eq:D-KE}\\
\widehat K_r^{\eta}F_i
 &=q^{-\delta_{ri}+\delta_{r,i+1}}
   F_i\widehat K_r^{\eta},
 \qquad \eta\in\{\pm\}.
 \label{eq:D-KF}
\end{align}

\item The mixed Chevalley relation is
\begin{equation}\label{eq:D-EF}
E_iF_j-F_jE_i
 =\delta_{ij}
 \frac{K_i^{+}-(K_i^{-})^{-1}}
      {q-q^{-1}}.
\end{equation}
\end{enumerate}

Moreover,
\begin{equation}\label{eq:D-factorization}
D_q\cong \Uc_q(\fb_-)\otimes \Uc_q(\fb_+)
\end{equation}
as a vector space.
\end{proposition}

\begin{proof}
This is the standard straightening calculation for a generalized quantum
double associated with a nondegenerate skew Hopf pairing.  The antipodes are
\[
S(E_i)=-(K_i^{+})^{-1}E_i,
\qquad
S(F_i)=-F_iK_i^{-},
\]
while all \(\widehat K_i^{\eta}\) are group-like.  Evaluation of the Majid
straightening formula on these group-like and skew-primitive generators gives
\eqref{eq:D-Cartan-commute}--\eqref{eq:D-EF}.  There are no mixed quantum
Serre relations because the double has the triangular factorization
\eqref{eq:D-factorization}.

For \(n=2\), the calculation is displayed explicitly in Proposition~2.10 and
Corollary~2.12 of Stokman's notes: after writing \(K\) and \(K_0\) for the two
Borel Cartans, one obtains
\[
K_0E=q^2EK_0,\qquad KF=q^{-2}FK,
\]
and
\[
EF-FE=\frac{K-K_0^{-1}}{q-q^{-1}},
\]
with \(\Uc_q(\mathfrak{sl}_2)\cong D/(K-K_0)\); see
\cite{StokmanDouble,StokmanSL2}.  The general-\(n\) formulas above are the
same calculation, with \(K_i^{\eta}\) replaced by the simple Cartan elements
\(\widehat K_i^{\eta}(\widehat K_{i+1}^{\eta})^{-1}\).  A related presentation
of quantum enveloping algebras by two Borel subalgebras is also given in
\cite{BB}; the Borel-double construction in rank one is treated in
\cite{BorelDouble}.
\end{proof}

\section{Root vectors and \(L\)-operators}

We next recall the standard root-vector normalization.  Put
\[
E_{i,i+1}:=E_i\qquad(1\le i<n)
\]
and, for \(i<j-1\), define inductively
\begin{equation}\label{eq:E-root}
E_{ij}:=E_{i,j-1}E_{j-1}
       -q^{-1}E_{j-1}E_{i,j-1}.
\end{equation}
Likewise put
\[
F_{i+1,i}:=F_i
\]
and, for \(i>j+1\),
\begin{equation}\label{eq:F-root}
F_{ij}:=F_{i-1,j}F_{i-1}
       -qF_{i-1}F_{i-1,j}.
\end{equation}
These are the standard \(q\)-root vectors attached to the positive and negative
roots
\[
\alpha_{ij}=\varepsilon_i-\varepsilon_j.
\]
The Borel PBW theorem says that ordered monomials in the \(E_{ij}\) and the
\(\widehat K_r^{+}\) form a basis of \(\Uc_q(\fb_+)\), while ordered monomials
in the \(F_{ij}\) and the \(\widehat K_r^{-}\) form a basis of
\(\Uc_q(\fb_-)\); see \cite[Chapter~6]{KS}.

\begin{definition}
In \(D_q\), define matrices \(L^+=(L^+_{ij})\) and \(L^-=(L^-_{ij})\) by
\begin{equation}\label{eq:Lplus}
L^+_{ij}=
\begin{cases}
0, & i>j,\\[1mm]
\widehat K_i^{+}, & i=j,\\[1mm]
\lambda E_{ij}\widehat K_j^{+}, & i<j,
\end{cases}
\end{equation}
and
\begin{equation}\label{eq:Lminus}
L^-_{ij}=
\begin{cases}
0, & i<j,\\[1mm]
(\widehat K_i^{-})^{-1}, & i=j,\\[1mm]
\lambda(\widehat K_j^{-})^{-1}F_{ij}, & i>j.
\end{cases}
\end{equation}
\end{definition}

For example, for \(n=2\),
\[
L^+=
\begin{pmatrix}
\widehat K_1^{+} & \lambda E_1\widehat K_2^{+}\\
0 & \widehat K_2^{+}
\end{pmatrix},
\qquad
L^-=
\begin{pmatrix}
(\widehat K_1^{-})^{-1} & 0\\
\lambda(\widehat K_1^{-})^{-1}F_1 & (\widehat K_2^{-})^{-1}
\end{pmatrix}.
\]

These matrices are the standard Faddeev--Reshetikhin--Takhtajan
\(L\)-operators, normalized so that the pairing between the simple Chevalley
generators is
\[
\la E_i,F_j\ra=\lambda^{-1}\delta_{ij}.
\]
In particular, their entries satisfy the usual triangular \(RLL\)-relations.
The factorization and PBW properties of these operators are standard in the
FRT description of \(\Uc_q^{\mathrm{ext}}(\sln)\); see \cite[Chapters~6 and~8]{KS}.

\section{The punctured-bigon presentation}

We now record the presentation of the stated skein algebra in the form most
convenient for comparison with \(D_q\).

\begin{theorem}[Punctured-bigon presentation]
\label{thm:punctured-bigon}
Let \(b_{ij}\) and \(c_{ij}\) be the horizontal right-directed arc generators
of \(\bcS_n(\fB^*)\).

\begin{enumerate}
\item The opposite triangular entries vanish:
\begin{equation}\label{eq:triangular}
b_{ij}=0\quad(i>j),
\qquad
c_{ij}=0\quad(i<j).
\end{equation}

\item The diagonal generators are invertible and satisfy
\begin{equation}\label{eq:det-one}
\prod_{i=1}^{n}b_{ii}=1,
\qquad
\prod_{i=1}^{n}c_{ii}=1.
\end{equation}

\item Define
\begin{equation}\label{eq:skein-ef}
e_i:=\lambda^{-1}b_{i,i+1}b_{i+1,i+1}^{-1},
\qquad
f_i:=\lambda^{-1}c_{ii}^{-1}c_{i+1,i}.
\end{equation}
Then \(b_{ii}\), \(c_{ii}\), \(e_i\), and \(f_i\) generate
\(\bcS_n(\fB^*)\), and their defining relations are exactly the relations of
\(D_q\) in Proposition~\ref{prop:double-presentation} under the dictionary
\[
b_{ii}\longleftrightarrow \widehat K_i^{+},
\qquad
c_{ii}^{-1}\longleftrightarrow \widehat K_i^{-},
\qquad
e_i\longleftrightarrow E_i,
\qquad
f_i\longleftrightarrow F_i.
\]

\item More generally, if \(E_{ij}\) and \(F_{ij}\) are defined recursively by
\eqref{eq:E-root} and \eqref{eq:F-root}, then in \(\bcS_n(\fB^*)\)
\begin{equation}\label{eq:skein-root-coordinates}
b_{ij}=\lambda E_{ij}b_{jj}\quad(i<j),
\qquad
c_{ij}=\lambda c_{jj}^{-1}F_{ij}\quad(i>j),
\end{equation}
where on the right-hand sides \(E_{ij}\) and \(F_{ij}\) mean the corresponding
quantum commutators of the elements \(e_r\) and \(f_r\) defined in
\eqref{eq:skein-ef}.

\item The algebra \(\bcS_n(\fB^*)\) has the PBW basis
\begin{equation}\label{eq:skein-PBW}
\mathbf F^{\,\mathbf r}\,
\mathbf c_{\mathrm{diag}}^{\,\mathbf a}\,
\mathbf b_{\mathrm{diag}}^{\,\mathbf d}\,
\mathbf E^{\,\mathbf m},
\end{equation}
where \(\mathbf E^{\,\mathbf m}\) and \(\mathbf F^{\,\mathbf r}\) are the usual
ordered monomials in the root vectors, the middle terms are Laurent monomials
in the diagonal arcs, and the product relations \eqref{eq:det-one} are imposed
on each diagonal family.
\end{enumerate}
\end{theorem}

\begin{proof}[Explanation and proof from the stated-skein calculus]
The bigon calculation of L\^e identifies the stated skein algebra of a bigon
with a quantum matrix algebra; splitting a punctured bigon along its standard
ideal arcs gives two such bigon factors.  The excision theorem for stated
skein algebras identifies the resulting gluing with the braided tensor product
of the two factors.  Written in terms of the horizontal arcs, this is exactly
the pair of triangular FRT matrices above.  These steps are contained in the
triangular-decomposition theorem of \cite{LeTri} and the excision and
punctured-surface calculations of \cite{CL}.

The triangular relations \eqref{eq:triangular}, determinant relations
\eqref{eq:det-one}, and PBW basis \eqref{eq:skein-PBW} are the corresponding
statements for the punctured bigon.  Substituting \eqref{eq:skein-ef} into the
quadratic skein relations gives the Borel relations for \(e_i\) and \(f_i\),
while the mixed relations between the two sheets give
\eqref{eq:D-Cartan-commute}--\eqref{eq:D-EF}.  The diamond-lemma part of the
presentation is equivalent to the stated PBW basis in
\eqref{eq:skein-PBW}.  Consequently there are no relations beyond those listed
in Proposition~\ref{prop:double-presentation}.
\end{proof}

\begin{remark}
The overall powers of \(q^{1/n}\) which often occur in stated
\(\mathrm{SL}_n\)-skein formulas have been absorbed into the normalization of
the horizontal arcs.  This is the normalization compatible with the pairing
\eqref{eq:pairing}.  Replacing the skein parameter by its inverse replaces
\(\lambda=q^{-1}-q\) by \(-\lambda\) throughout.
\end{remark}

\section{The explicit isomorphism}

\begin{theorem}\label{thm:main-isomorphism}
There is a unique algebra isomorphism
\[
\Phi:\bcS_n(\fB^*)\xrightarrow{\;\sim\;}
\Uc_q(\fb_-)\bipr\Uc_q(\fb_+)
\]
such that
\begin{equation}\label{eq:Phi-b}
\Phi(b_{ij})=
\begin{cases}
0, & i>j,\\
\widehat K_i^{+}, & i=j,\\
\lambda E_{ij}\widehat K_j^{+}, & i<j,
\end{cases}
\end{equation}
and
\begin{equation}\label{eq:Phi-c}
\Phi(c_{ij})=
\begin{cases}
0, & i<j,\\
(\widehat K_i^{-})^{-1}, & i=j,\\
\lambda(\widehat K_j^{-})^{-1}F_{ij}, & i>j.
\end{cases}
\end{equation}
Its inverse is determined by
\begin{equation}\label{eq:Psi-K}
\Phi^{-1}(\widehat K_i^{+})=b_{ii},
\qquad
\Phi^{-1}(\widehat K_i^{-})=c_{ii}^{-1},
\end{equation}
and
\begin{equation}\label{eq:Psi-EF}
\Phi^{-1}(E_i)
 =\lambda^{-1}b_{i,i+1}b_{i+1,i+1}^{-1},
\qquad
\Phi^{-1}(F_i)
 =\lambda^{-1}c_{ii}^{-1}c_{i+1,i}.
\end{equation}
\end{theorem}

\begin{proof}
By \eqref{eq:triangular}, formulas \eqref{eq:Phi-b} and \eqref{eq:Phi-c} are
consistent with the triangularity relations.  By Theorem
\ref{thm:punctured-bigon}(3), the assignment
\[
b_{ii}\longmapsto \widehat K_i^{+},
\qquad
c_{ii}\longmapsto (\widehat K_i^{-})^{-1},
\qquad
e_i\longmapsto E_i,
\qquad
f_i\longmapsto F_i
\]
extends to a homomorphism.  Formula \eqref{eq:skein-root-coordinates} shows
that the same homomorphism sends the arbitrary horizontal generators to
\eqref{eq:Phi-b} and \eqref{eq:Phi-c}.

The inverse formulas are forced and are compatible with the relations.  For
example,
\[
\begin{aligned}
\Phi\bigl(\lambda^{-1}b_{i,i+1}b_{i+1,i+1}^{-1}\bigr)
&=\lambda^{-1}
 \bigl(\lambda E_i\widehat K_{i+1}^{+}\bigr)
 (\widehat K_{i+1}^{+})^{-1}\\
&=E_i,
\end{aligned}
\]
and
\[
\begin{aligned}
\Phi\bigl(\lambda^{-1}c_{ii}^{-1}c_{i+1,i}\bigr)
&=\lambda^{-1}\widehat K_i^{-}
 \bigl(\lambda(\widehat K_i^{-})^{-1}F_i\bigr)\\
&=F_i.
\end{aligned}
\]
Similarly,
\[
\Phi^{-1}\Phi(b_{ii})=b_{ii},
\qquad
\Phi^{-1}\Phi(c_{ii})=c_{ii},
\]
and the corresponding identities on the off-diagonal arcs follow from
\eqref{eq:E-root}, \eqref{eq:F-root}, and
\eqref{eq:skein-root-coordinates}.

Finally, the PBW basis \eqref{eq:skein-PBW} of \(\bcS_n(\fB^*)\) is mapped to
the tensor-factor PBW basis of \(D_q\) in
\eqref{eq:D-factorization}: a positive-root monomial is sent to the
corresponding monomial in the \(E_{ij}\), a negative-root monomial to the
corresponding monomial in the \(F_{ij}\), and the two diagonal Laurent
monomials to the two Cartan Laurent monomials.  The scalar power of
\(\lambda\) is the number of off-diagonal factors and is nonzero.  Hence
\(\Phi\) sends a basis to a basis, so it is bijective.
\end{proof}

\section{The quotient by \(b_{ii}c_{ii}-1\)}

Let \(I\subset\bcS_n(\fB^*)\) be the two-sided ideal
\begin{equation}\label{eq:I}
I:=\bigl(b_{ii}c_{ii}-1\mid 1\le i\le n\bigr).
\end{equation}
Likewise, let \(J\subset D_q\) be
\begin{equation}\label{eq:J}
J:=\left(\widehat K_i^{+}(\widehat K_i^{-})^{-1}-1
       \mid 1\le i\le n\right).
\end{equation}

\begin{lemma}\label{lem:J-central}
The elements
\[
g_i:=\widehat K_i^{+}(\widehat K_i^{-})^{-1}
\]
are central group-like elements of \(D_q\).  Consequently \(J\) is a Hopf
ideal.
\end{lemma}

\begin{proof}
Both factors are group-like, so \(g_i\) is group-like.  By
\eqref{eq:D-KE},
\[
\widehat K_i^{+}E_j
=q^{\delta_{ij}-\delta_{i,j+1}}E_j\widehat K_i^{+},
\]
whereas
\[
(\widehat K_i^{-})^{-1}E_j
=q^{-\delta_{ij}+\delta_{i,j+1}}
 E_j(\widehat K_i^{-})^{-1}.
\]
Thus \(g_i\) commutes with \(E_j\).  The same argument, with the inverse
exponents, shows that \(g_i\) commutes with \(F_j\).  It clearly commutes with
both Cartan tori.  Therefore \(g_i\) is central.  Moreover
\[
\varepsilon(g_i)=1,
\qquad
S(g_i)=g_i^{-1},
\]
so the ideal generated by \(g_i-1\) is a Hopf ideal.
\end{proof}

\begin{theorem}\label{thm:quotient}
Under the isomorphism \(\Phi\),
\[
\Phi(I)=J.
\]
Consequently there is a canonical algebra isomorphism
\begin{equation}\label{eq:main-quotient}
\bcS_n(\fB^*)/I
\;\cong\;
\Uc_q^{\mathrm{ext}}(\sln).
\end{equation}
The isomorphism sends
\[
[b_{ii}]\longmapsto\widehat K_i,
\qquad
[c_{ii}]\longmapsto\widehat K_i^{-1},
\]
and
\[
[\lambda^{-1}b_{i,i+1}b_{i+1,i+1}^{-1}]
 \longmapsto E_i,
\qquad
[\lambda^{-1}c_{ii}^{-1}c_{i+1,i}]
 \longmapsto F_i.
\]
\end{theorem}

\begin{proof}
Since Cartan elements from the two Borel factors commute in \(D_q\),
\[
\Phi(b_{ii}c_{ii})
=\widehat K_i^{+}(\widehat K_i^{-})^{-1}
=g_i.
\]
Thus \(\Phi(I)=J\).

Now define
\[
\pi:D_q\longrightarrow \Uc_q^{\mathrm{ext}}(\sln)
\]
by
\[
\pi(E_i)=E_i,\qquad
\pi(F_i)=F_i,\qquad
\pi(\widehat K_i^{+})=\widehat K_i,\qquad
\pi(\widehat K_i^{-})=\widehat K_i.
\]
This preserves all relations in
Proposition~\ref{prop:double-presentation}.  The only relation requiring a
check is \eqref{eq:D-EF}; its image is
\[
[E_i,F_j]
=\delta_{ij}
\frac{K_i-K_i^{-1}}{q-q^{-1}},
\]
because
\[
\pi(K_i^{+})=K_i,
\qquad
\pi((K_i^{-})^{-1})=K_i^{-1}.
\]
Hence \(\pi\) is well-defined.  Clearly \(J\subseteq\ker\pi\).

Conversely, define
\[
s:\Uc_q^{\mathrm{ext}}(\sln)\longrightarrow D_q/J
\]
by sending \(E_i,F_i,\widehat K_i\) to the classes of
\(E_i,F_i,\widehat K_i^{+}\).  In \(D_q/J\),
\[
[\widehat K_i^{+}]=[\widehat K_i^{-}],
\]
so every defining relation of \(\Uc_q^{\mathrm{ext}}(\sln)\) is preserved.
Hence \(s\) is well-defined.  The composite
\[
\Uc_q^{\mathrm{ext}}(\sln)
 \xrightarrow{s}D_q/J
 \xrightarrow{\bar\pi}\Uc_q^{\mathrm{ext}}(\sln)
\]
is the identity.  Also, \(D_q/J\) is generated by the classes in the image of
\(s\), so \(s\) is surjective.  Therefore \(s\) and \(\bar\pi\) are mutually
inverse, proving
\[
D_q/J\cong\Uc_q^{\mathrm{ext}}(\sln).
\]
Applying \(\Phi^{-1}\) gives \eqref{eq:main-quotient}.  The displayed images of
the quotient generators are immediate from
\eqref{eq:Psi-K}--\eqref{eq:Psi-EF}.
\end{proof}

\begin{corollary}
The quotient \(\bcS_n(\fB^*)/I\) is a Hopf algebra, and the isomorphism in
Theorem~\ref{thm:quotient} is an isomorphism of Hopf algebras.
\end{corollary}

\begin{proof}
The generalized double is a Hopf algebra, and \(J\) is a Hopf ideal by
Lemma~\ref{lem:J-central}.  The quotient map identifies the two Borel Cartan
tori.  Its coproduct therefore becomes
\[
\Delta(E_i)=E_i\otimes1+K_i\otimes E_i,
\qquad
\Delta(F_i)=F_i\otimes K_i^{-1}+1\otimes F_i,
\]
and \(\Delta(\widehat K_i)=\widehat K_i\otimes\widehat K_i\), which are the
coproducts in \(\Uc_q^{\mathrm{ext}}(\sln)\).
\end{proof}

\begin{remark}
For \(n=2\), Theorem~\ref{thm:quotient} is precisely the relation
\[
D/(K-K_0)\cong \Uc_q(\mathfrak{sl}_2)
\]
displayed in Corollary~2.12 of \cite{StokmanDouble}.  In that notation, the
two diagonal products \(b_{11}c_{11}\) and \(b_{22}c_{22}\) give one
independent relation because \(b_{11}b_{22}=c_{11}c_{22}=1\).  The same
dependence occurs for general \(n\): although \(I\) is written with \(n\)
generators, \eqref{eq:det-one} makes their product redundant.
\end{remark}

\begin{thebibliography}{99}

\bibitem[KS]{KS}
A.~Klimyk and K.~Schm\"udgen,
\emph{Quantum Groups and Their Representations},
Texts and Monographs in Physics, Springer, 1997.
\newblock \href{https://link.springer.com/book/10.1007/978-3-642-60834-6}{link.springer.com}.

\bibitem[Le]{LeTri}
T.~T.~Q. L\^e,
\emph{Triangular decomposition of skein algebras},
Quantum Topol. \textbf{9} (2018), 591--632.
\newblock \href{https://arxiv.org/abs/1609.04987}{arxiv.org}.

\bibitem[CL]{CL}
F.~Costantino and T.~T.~Q. L\^e,
\emph{Stated skein algebras of surfaces},
J. Eur. Math. Soc. \textbf{25} (2023), 4063--4142.
\newblock \href{https://arxiv.org/abs/1907.11400}{arxiv.org}.

\bibitem[St1]{StokmanDouble}
J.~V. Stokman,
\emph{Quantum double as a braided Hopf algebra and quantum \(\mathfrak{sl}_2\)}.
\newblock \href{https://staff.fnwi.uva.nl/j.v.stokman/quantumdouble2.pdf}{staff.fnwi.uva.nl}.

\bibitem[St2]{StokmanSL2}
J.~V. Stokman,
\emph{Towards a universal \(R\)-matrix for \(\Uc_q(\mathfrak{sl}_2)\)}.
\newblock \href{https://staff.fnwi.uva.nl/j.v.stokman/Quantumsl2.pdf}{staff.fnwi.uva.nl}.

\bibitem[BB]{BB}
A.~Berenstein and S.~Braverman,
\emph{Presentation by Borel subalgebras and Chevalley generators for quantum
enveloping algebras},
Proc. Edinburgh Math. Soc.
\newblock \href{https://www.cambridge.org/core/services/aop-cambridge-core/content/view/886B0B4D8E3D242A51A0380A486351EA/S0013091504000689a.pdf/presentation-by-borel-subalgebras-and-chevalley-generators-for-quantum-enveloping-algebras.pdf}{cambridge.org}.

\bibitem[BD]{BorelDouble}
\emph{The quantum double of the Borel subalgebra of \(\Uc_q(\mathfrak{sl}_2)\)}.
\newblock \href{https://arxiv.org/pdf/math/0512563}{arxiv.org}.

\end{thebibliography}

\end{document}
```